\chapter{Двухрезервуарный неравновесный ток
\texorpdfstring{$K_{43}$}{K43}}
\label{ch:v10-k43-two-reservoir-current}

Возьмём две ванны с общей нисходящей вероятностью на канал
$p_\downarrow=1/6$, но различными KMS-отношениями:
\begin{equation}
 r_h=\frac12=e^{-\log2},\qquad r_c=\frac14=e^{-2\log2}.
\end{equation}
Тогда восходящие вероятности равны $1/12$ и $1/24$. Суммарный CPTP-канал
имеет матрицу популяций
\begin{equation}
 T=\begin{pmatrix}7/8&1/3\\1/8&2/3\end{pmatrix},
 \qquad \pi_*=(8/11,3/11)^T.
\end{equation}
Частичные стационарные токи равны
\begin{equation}
 J_h=\frac1{12}\frac8{11}-\frac16\frac3{11}=\frac1{66},
 \qquad
 J_c=\frac1{24}\frac8{11}-\frac16\frac3{11}=-\frac1{66}.
\end{equation}
Система не накапливает населённость, но через неё проходит ненулевой ток
$1/66$ от горячей ванны к холодной. Производство энтропии положительно:
\begin{equation}
 \dot S=J_h(2\log2-\log2)=\frac{\log2}{66}>0.
\end{equation}

\begin{theorem}[Условный двухрезервуарный ток]
Две KMS-ванны с различными сродствами дают точный CPTP-канал,
стационарный сквозной ток $1/66$ и положительное производство энтропии.
\end{theorem}

\begin{nogocriterion}[Архитектура не выводит ванны]
Значения двух сродств и общий дискретный темп заданы условно. Канал
доказывает возможность неравновесного дыхания, но не физическое происхождение
резервуаров, абсолютной скорости или масштаба.
\end{nogocriterion}

\begin{protocol}\begin{enumerate}
\item CPTP-канал и стационарность: \textbf{получены};
\item ненулевой ток: \textbf{$1/66$};
\item производство энтропии: \textbf{$\log2/66>0$};
\item условная архитектура: \textbf{$10/10$};
\item происхождение ванн, сродств и темпа: \textbf{$0/3$};
\item абсолютный масштаб: \textbf{$0/1$};
\item следующий гейт: \textbf{хопфовское происхождение разности сродств}.
\end{enumerate}\end{protocol}

Машинный сертификат сохранён в
\path{s2t_v10_cell_birth_four_volume_induced_newton_breathing_anomaly_k43_nonequilibrium_two_reservoir_output_current_parent_admission_gate_results.json}.